\chapter{Protected Compact Cores and Prime Replacement}
\label{ch:compact-cores}

Two constructions make the lower-handle argument relative to previously
fixed geometry. The first encloses a prescribed compact set and the entire
old boundary in a compact PL domain with just one new spherical boundary.
The second changes the PL structure on a lattice handle to an irreducible
one while retaining the original structure near a specified ball and the
entire old boundary. An abstract classification of the resulting manifold
would not supply these protection conditions.

\section{Domains, parametrizations, and ends}
\label{sec:compact-core-definitions}

Write $V^d=\mathbb R^d$ with the maximum norm, $B^d=[-1,1]^d$, and
$Q^{d-1}=\partial B^d$. These are the literal ball and sphere models used
here. PL parametrizations below have cubical source boundaries. Let $X$
be Hausdorff and let $e=(e_i)$ be a family of open partial homeomorphisms
from $X$ to $V^3$. A map is chartwise PL on a set when its coordinate
restrictions have local finite polyhedral, piecewise affine descriptions.

\begin{definition}[Compatible PL domain]
\label{def:compact-pl-domain}
A PL domain $R$ for $e$ is closed in $X$; the charts $e_i$ cover $X$ and
have piecewise affine transitions. For every $x\in\operatorname{fr}_X R$
there are a compatible partial chart $b$, an affine functional $\ell$,
and $v\in V^3$ such that
\[
 x\in\operatorname{source}(b),\qquad
 \ell(b(x))=0,\qquad \ell_{\rm lin}(v)=1,
\]
\[
 y\in R\quad\Longleftrightarrow\quad\ell(b(y))\geq0
 \qquad(y\in\operatorname{source}(b)).
\]
Compatibility requires every transition $b\circ e_i^{-1}$ on its actual
domain to be piecewise affine. The condition on $\ell_{\rm lin}$ makes
the half-space nondegenerate.
\end{definition}

Let $j:R\hookrightarrow X$ denote inclusion. A statement about
$\operatorname{int}_R j^{-1}K$ differs from one about
$\operatorname{int}_X K$: the whole old boundary can belong to the former
while belonging to $\operatorname{fr}_X K$.

\begin{definition}[One simply connected end]
\label{def:one-simply-connected-end}
A space $Y$ has one simply connected end when for every compact
$C\subset Y$ there exists compact $D\subset Y$ such that
$C\subset\operatorname{int}_Y D$, $Y\setminus D$ is nonempty and connected,
and every based loop $p:[0,1]\to Y$ whose image avoids $D$ admits a
based nullhomotopy in $Y\setminus C$.
\end{definition}

The compact set $D$ depends on $C$, and every loop outside this same $D$
must contract while avoiding $C$. The contraction may meet $D\setminus C$;
simple connectedness of $Y\setminus D$ is not assumed. The condition
implies that $Y$ is noncompact: applying it to $C=Y$ would require a
nonempty complement of a set containing $Y$.

\begin{definition}[Marked PL sphere, ball, and irreducibility]
\label{def:chartwise-sphere-ball}
A chartwise PL sphere on $S\subset X$ consists of a homeomorphism
$s:Q^2\to S$ and a map $\widetilde s:V^3\to X$ equal to $s$ on $Q^2$
and chartwise PL there. No regularity away from $Q^2$ is required.
A chartwise PL ball on $(D,S)$ consists of $S\subset D$, a homeomorphism
$b:B^3\to D$, and a map $V^3\to X$ agreeing with $b$ and chartwise PL
on $B^3$, such that $b(u)\in S$ if and only if $u\in Q^2$.
A PL domain $R$ is irreducible in $e$ when every such sphere
$S\subset\operatorname{int}_X R$ bounds such a ball $(D,S)$ with $D\subset R$.
\end{definition}

The boundary set and the entire parametrization are retained. A topological
ball filling by itself does not establish this PL conclusion.

\section{Protected domains from the weak end}
\label{sec:weak-end-core}

Assume $R$ is a connected PL domain with compact old frontier
$B=\operatorname{fr}_X R$, that $R$ has one simply connected end, and that
$A\subset R$ is compact. First construct compact connected PL domains
$C\subset K\subset R$ with
\[
 j^{-1}(A\cup B)\subset\operatorname{int}_R j^{-1}C,
 \qquad j^{-1}C\subset\operatorname{int}_R j^{-1}K.
\]
The domain $C$ will stay fixed during compression. We require a nonempty
compact $F\subset\operatorname{int}_X R$ with
\begin{equation}
 \operatorname{fr}_X K=B\cup F,\qquad
 \operatorname{fr}_R j^{-1}K=j^{-1}F,\qquad B\cap F=\varnothing,
 \label{eq:protected-frontiers}
\end{equation}
and every loop on $F$, viewed in $R$, must contract in $R\setminus C$.

The compact PL neighborhood construction first gives a connected $C$
protecting $A\cup B$. Finite chart neighborhoods around this compact set
are joined by finitely many PL paths in the connected domain before
thickening; half-space charts keep the old boundary on the domain side.
Apply the end condition to $j^{-1}C$ and obtain $D\subset R$.
Choose a compact ambient PL neighborhood $P$ with
$C\cup B\cup j(D)\subset\operatorname{int}_X P$ and set
\[
 K_0=P\cap R,\qquad F_0=\operatorname{fr}_X P\cap R.
\]
The boundaries of $P$ and $R$ are separated because $B$ lies in the
interior of $P$. Their half-space charts give a PL structure on $K_0$
and the exact equalities
\[
 \operatorname{fr}_X K_0=B\cup F_0,
 \qquad\operatorname{fr}_R j^{-1}K_0=j^{-1}F_0.
\]
Moreover $F_0\cap j(D)=\varnothing$. Hence the chosen end contraction
applies to every loop in $F_0$ and avoids $C$.

Take the component $K$ of $K_0$ containing $C$. PL half-space domains are
locally path connected, so this component is relatively open in $K_0$;
it is also closed and compact. Its new frontier is $F=F_0\cap K$.
Relative openness preserves the domain charts and both frontier identities.
Connectedness puts all of $C$ in this one component, including the old
boundary. The new frontier is nonempty: otherwise $j^{-1}K$ is nonempty
and clopen in connected $R$, and therefore equals $R$, contradicting its
noncompactness. Every loop in $F$ is a loop in $F_0$, so the same
nullhomotopy works pointwise.

Here the half-space assertion is literal. The compact set
$A\cup B\cup j(D)$ is covered by finitely many actual chart neighborhoods.
After shrinking and taking a common polyhedral refinement, join the finitely
many pieces meeting $A\cup B$ by paths in the connected domain and thicken
the resulting finite carrier. At a point of $B$ the old boundary chart is
restricted to the interior of $P$; at a point of $F_0$ the chart for $P$ is
restricted to the interior of $R$. Since $B\subset\operatorname{int}_X P$,
these two kinds of charts do not meet in a corner. Thus
\[
 \operatorname{fr}_X(P\cap R)=B\cup(\operatorname{fr}_X P\cap R),
 \qquad
 \operatorname{fr}_R j^{-1}(P\cap R)=j^{-1}(\operatorname{fr}_X P\cap R).
\]
The component of $P\cap R$ containing the connected carrier is relatively
open because half-space domains are locally path connected. It is therefore
clopen in $P\cap R$, so restricting the same charts preserves both frontier
identities. This is the finite polyhedral construction in the
protected-neighborhood and connected-core producers; ordinary compactness
alone would not give these identities.

\section{Compression and its decreasing measure}
\label{sec:compression-minimality}

Represent $F$ by a finite two-dimensional simplicial complex with actual
component images $S_1,\ldots,S_n$. The same embedding and component inverse
maps are retained throughout. The genus labels obey
\[
 v_i+t_i+2g_i=e_i+2,
\]
where $v_i,e_i,t_i$ count vertices, edges, and triangles of the component.
The orientability and genus construction uses the protected loop
contractions; ambient orientability is not an extra hypothesis. The
finite surface analysis also gives a PL sphere parametrization for $g_i=0$.

\begin{definition}[Compression complexity]
For this finite component model put
\[
 w(g)=\max\{2g-1,0\},\qquad c(F)=\sum_{i=1}^n w(g_i).
\]
No vertex, edge, or triangle count is added to this measure.
\end{definition}

Fix $C$. An admissible witness consists of $(K,F)$ satisfying all the
conditions above, including connectedness, both frontier identities, the
loop contractions avoiding $C$, and a realized finite component model
with the Euler identities. The preceding construction and finite surface
analysis give one such witness. The nonempty set of realized natural-number
complexities has a least element. Choose a witness attaining it.

\begin{proposition}[No essential proper disk for a minimal core]
\label{prop:compact-core-incompressibility}
For this witness, let $d:B^2\to R\setminus C$ be an embedded chartwise PL
disk such that $d^{-1}(F)=Q^1$. Its boundary loop, with the complete
$Q^1$ parametrization, represents the identity in the fundamental group
of the component of $F$ containing it.
\end{proposition}
\begin{proof}
Suppose the loop is essential. Connectedness puts the disk interior on
one side of $F$. Compress in a product neighborhood of the disk, replacing
the boundary annulus by the two disk faces. On the interior side this
cuts the core; on the exterior side it enlarges it. All changes lie in
$R\setminus C$. The protected compression construction retains the
component containing $C$, supplies its half-space charts, and transports
the frontier-loop contractions through the cut. It is thus admissible
for the same $C$.

A nonseparating compression replaces positive genus $g$ by $g-1$.
For $g=1$ the weight drops from $1$ to $0$; for $g\geq2$ it drops by $2$.
A separating essential compression produces positive genera $a,b$ with
$g=a+b$, and
\[
 (2a-1)+(2b-1)=2g-2<2g-1.
\]
The positivity assumptions are necessary: a genus-zero side would make
the original rim inessential. The finite cap-piece Euler count gives
these identities for the actual successor model. Discarding components
outside the protected one cannot increase the sum of nonnegative weights.
The successor therefore has strictly smaller complexity, a contradiction.
\end{proof}

For clarity, the product operation is relative to the protected set. Put
$Y=R\setminus C$. Because $B\subset C$ and $R$ is closed,
$Y=\operatorname{int}_X R\setminus C$ is open. A complete bicollar
$h:F\times(-1,1)\to U\subset Y$ satisfies
$h(x,t)\in K$ exactly when $t\geq0$ and $h(x,t)\in F$ exactly when
$t=0$, simultaneously for every component of $F$. The image
$d(\operatorname{int}B^2)$ is connected and avoids $F$, so it lies in one
of the two open sides of $Y\setminus F$.
In its selected closed side construct an embedded chartwise PL disk product
$P:B^2\times[-1,1]\to Y$ with
\[
 P(z,0)=d(z),\qquad P(z,s)\in F\quad\Longleftrightarrow\quad z\in Q^1.
\]
For $0<\epsilon<1$, the lateral annulus
$P(Q^1\times[-\epsilon,\epsilon])$ lies in $F$. Remove its open strip
$P(Q^1\times(-\epsilon,\epsilon))$ and insert the two disk faces
$P(B^2\times\{-\epsilon,\epsilon\})$. In the interior case this removes
the open disk-product strip from $K$; in the exterior case it adjoins the
closed strip. The disk-product embedding and its exact frontier preimage
give the embedded replacement and its complete rim identifications.
The product lies in $Y$, so the operation misses $C$. The finite cap-piece
count gives the Euler identities quoted above; the component containing $C$ is then
selected before discarding any other components. Thus the successor retains
the exact relative interior and loop-contraction hypotheses rather than
merely producing an abstract compressed surface.

Total genus would not suffice, since separating compression preserves
$a+b$. The proposition concerns boundary loops of these proper disks;
it does not assert that all frontier loops are trivial.

\section{The spherical alternative}
\label{sec:spherical-frontier}

The marked-boundary Dehn theorem takes a PL domain, an open marked part
of its boundary, and a continuous disk filling of a loop whose class is
nontrivial in that marked part. It returns an embedded proper PL disk
with essential boundary there. The class, marking, and filling must
refer to the same cut domain. This is an input from the Dehn and surfaces
chapter.

Form the protected region $Y=R\setminus C$. It is open in $X$, because
$C$ contains $\operatorname{fr}_X R$, and it contains $F$. The finite
surface construction gives an alternative for each component of $F$:
either it is a parametrized PL sphere, or it supplies an essential marked
class with a continuous filling in a cut domain inside $Y$. The end
condition supplied the initial nullhomotopy. Cutting and lifting that
nullhomotopy supplies the one-sided filling; a general nullhomotopy in
$Y$ does not automatically lie on the required side.

In the second case the marked part is the complete boundary of the cut
domain, hence open relative to that boundary. Apply the marked-boundary
theorem with the trivial subgroup. Projection through the cut charts
gives an embedded chartwise PL disk in $R\setminus C$, with interior off
$F$ and an essential boundary class in $F$. This contradicts
Proposition~\ref{prop:compact-core-incompressibility}. Thus the new
frontier is a finite pairwise disjoint family of actual PL spheres in
$\operatorname{int}_X R$.

Two finite constructions are needed for this alternative. First, the
weak-end homotopies orient each component of the frontier: orient one
triangle, transport the sign across a spanning tree of its dual graph, and
use the nullhomotopy of every dual loop to show that transport is independent
of the chosen path. These signs come from the ambient orientation double
cover restricted to the frontier, with its fixed outward coorientation;
the ambient nullhomotopies therefore make its monodromy trivial.
In the finite tree-cotree calculation, the copied polygon
initially has two signed copies of every primal-tree edge as well as of every
residual edge. The primal-tree contraction is performed after the dual-tree
gluing; only then may the boundary word be indexed by residual edges. Dropping
the primal-tree pairs at the first stage gives the tetrahedral counterexample
recorded in the correction note and would make the orientation count false.

Second, the nullhomotopy is made one-sided before Dehn is invoked. Inside
$Y$ use the two closed domains $K_+=K\cap Y$ and
$K_-=Y\setminus\operatorname{int}_X K$ in the same lifted chart family.
Their complete frontier is $F$, their interiors are disjoint, and their union
is $Y$. The balanced square-rim quotient turns the based nullhomotopy into a
disk with exactly the prescribed rim values. A finite height perturbation in
the bicollar makes the preimage of $F$ polygonal. Retain the collar and
approximation homotopies from the original rim to the perturbed rim in
$\operatorname{int}_Y K_+$. Delete innermost inessential circles while
fixing that perturbed rim. If an innermost essential circle is reached,
its subdisk gives a proper filling in one of the two closed sides; its
essential class lies in the whole $F$ and may belong to a different
component from the initial rim. If every preimage circle disappears, the
remaining disk avoids $F$ altogether. Its connected image and its boundary
in $\operatorname{int}_Y K_+$ place it entirely in that interior. Paste
this disk to the two retained rim homotopies in successive radial annuli.
The resulting filling restores the original essential rim pointwise and
has interior disjoint from $F$. Thus both alternatives give an essential
proper filling on a single side. The marked Dehn producer then supplies
the proper embedded PL disk used in the contradiction. The finite approximation and
prescribed-boundary producer remain the explicit input from the Dehn-surfaces
chapter; no unmarked Dehn lemma is substituted.

The cut-filling alternative and disk projection are distinct geometric
steps. The surface and covering arguments above are needed in addition
to incompressibility to recognize each frontier component as a sphere.

\section{Filling the bounded complementary components}
\label{sec:compact-core-filling}

Only the connected-tail part of the end condition is needed here. Put
$k=j^{-1}K$. Choose compact $d\subset R$ with
$k\subset\operatorname{int}_R d$ and connected nonempty complement.
For $x\in R\setminus d$, let $U$ be its component in $R\setminus k$
and set $l=R\setminus U$. The entire tail $R\setminus d$ lies in $U$,
so $l\subset d$. Local connectedness makes $U$ open; hence $l$ is
closed and compact, and $k\subset l$.

Every other complementary component $E$ has closure in $l$. Its frontier
lies in $k$, since a point outside $k$ has a neighborhood in its own open
component. That frontier is nonempty by connectedness of $R$ and
nonemptiness of $k$. Therefore $k\cup\overline E$ is connected.
The union of these sets and $k$ is exactly $l$, proving connectedness of
$l$ without assuming finitely many complementary components.

The full bicollar of $F=\bigcup_iS_i$ determines which spheres remain.
Choose positive height on the $K$ side. Each connected negative half-collar
of $S_i$ lies in one complementary component $E_i$ and accumulates at every
point of $S_i$. Retain exactly $J=\{i:E_i=U\}$. Then
\[
 \operatorname{fr}_R U=j^{-1}\left(\bigcup_{i\in J}S_i\right).
\]
Indeed, a boundary point of $U$ lies on $F$; its open bicollar neighborhood
meets $U$ on the negative side, forcing $E_i=U$. Conversely the entire
base sphere is approached from that half-collar. Whole spheres are
selected, rather than fragments.

Set $L=j(l)$. Near each retained sphere $L$ agrees with $K$, inheriting
the same half-space charts; near the old boundary the original relative
neighborhood remains. Elsewhere there is no frontier. Thus $L$ is a
compact connected PL domain and
\[
 K\subset L\subset R,\qquad
 \operatorname{fr}_X L=B\cup\bigcup_{i\in J}S_i,
 \qquad\operatorname{fr}_R j^{-1}L=j^{-1}\bigcup_{i\in J}S_i.
\]
Its relative complement is connected and nonempty. The index set $J$ is
nonempty, since otherwise $l$ is a nonempty compact clopen subset of
noncompact connected $R$. Inclusion $K\subset L$ preserves every protected
relative interior. No filled complementary component is asserted to be a ball.

\section{Joining two boundary spheres}
\label{sec:sphere-merging}

Suppose a filled core $L$ has $n>1$ new spheres and protect the closed set
$P=A\cup B$. Choose points on distinct spheres $S_a,S_b$. Since
$B\subset L$, the connected exterior $R\setminus L$ equals the open set
$\operatorname{int}_X R\setminus L$. Half-space charts give outward exits
from the chosen points. A PL path joins the exits in this exterior, and
a simple arc in its finite polyhedral image removes repeated excursions.
It yields an embedded PL arc $q$ whose only contacts with $L$ are its
endpoints and whose whole image lies in $\operatorname{int}_X R\setminus P$.

Thicken $q$ to a PL ball with two feet, one on each sphere. The construction
uses finite coordinate tubes along a subdivision of $q$ and half-space
charts at its endpoints. Choose a second point on each sphere and exclude
it from the tube neighborhood, so that neither foot consumes its entire
sphere. The tubes avoid the other spheres and $P$ and are pasted with
their full product identifications. The resulting ball meets $L$ in
exactly its two disk feet. The new boundary is the two punctured spheres
joined by the lateral annulus, hence a PL sphere $T$: the two residual
disks cap the boundary circles of the annulus, with the complete rim
parametrizations used in the pasting.

The larger core $K'$ satisfies
\[
 \operatorname{fr}_X K'=B\cup T\cup\bigcup_{i\ne a,b}S_i,
 \qquad K'\setminus L\subset\operatorname{int}_X R\setminus P.
\]
It agrees with $L$ on an open neighborhood of $P$. The sphere $T$ lies
in $\operatorname{int}_X R$, since it is disjoint from the retained old
boundary. The ambient frontier identity and protection imply the relative
one: ambient and relative interiors agree at interior points of $R$,
while the old boundary is in the relative interior of the core.

The intermediate family has $n-1$ members. Fill its bounded complementary
components again, retaining a nonempty subfamily of size $m\leq n-1$.
Strong induction on the positive sphere count ends at one sphere. Its
invariant contains compactness, connectedness, both exact frontier
identities, inclusion in the same $R$, and the protected relative interior.
All cores increase, so the original compact set stays protected.

\begin{theorem}[Protected Wall core]
\label{thm:protected-wall-core}
Let $X$ be Hausdorff and paracompact, $e$ an ambient chart family, and $R$
a connected PL domain with compact frontier and one simply connected end.
For every compact $A\subset R$ there are compact $K\subset R$ and a
chartwise PL sphere $S\subset\operatorname{int}_X R$ such that $K$ is a
PL domain and
\[
 \operatorname{fr}_X K=\operatorname{fr}_X R\cup S,\qquad
 S\cap\operatorname{fr}_X R=\varnothing,
\]
\[
 j^{-1}(A\cup\operatorname{fr}_X R)\subset\operatorname{int}_R j^{-1}K,
 \qquad\operatorname{fr}_R j^{-1}K=j^{-1}S.
\]
The construction also chooses $K$ connected.
\end{theorem}
\begin{proof}
The minimal protected core and marked-disk contradiction give a finite
spherical frontier. Apply the filling construction and sphere induction
with $P=A\cup B$. Compactness of $S$ follows from its parametrization.
\end{proof}

\section{The marked lattice-handle contract}
\label{sec:marked-handles}

Let $I,J$ be finite sets with cardinalities $i,k$ satisfying $i+k=3$ and
$i\leq2$. Let $\Lambda\subset V^J$ be a discrete full integer lattice:
it has a $\mathbb Z$-basis which is also a real basis of $V^J$. Use the
literal ambient space, domain, and quotient projection
\[
 X=V^I\times(V^J/\Lambda),\qquad
 R=B^I\times(V^J/\Lambda),\qquad \pi(x,y)=(x,[y]).
\]
Writing $B_r^J=\{y:\|y\|_\infty\leq r\}$, set
$C_r=\pi(B^I\times B_r^J)$ and
$A_r=\pi(\partial B^I\times B_r^J)$.

\begin{definition}[Protected marked ball]
For an ambient PL chart family $e$ making $R$ a PL domain, a protected
marked ball is a chartwise PL ball $(D,\operatorname{fr}_X D)$ with
$D\subset R$ and these position conditions. If $i=0$, require
$D\subset\operatorname{int}_X R$. If $i=1$ or $2$, require
\[
 C_1\subset D\subset C_2,\qquad
 D\subset\pi(B^I\times\{y:\|y\|_\infty<2\}),
\]
\[
 j^{-1}C_1\subset\operatorname{int}_R j^{-1}D,\qquad
 D\cap\operatorname{fr}_X R=A_{3/2}.
\]
\end{definition}

Existence of this ball for every lattice is not asserted; it is a
conditional input to replacement. The lower-handle application constructs
the ball after choosing its periods.

\begin{theorem}[Protected irreducible replacement]
\label{thm:protected-prime-replacement}
For every $I,J,\Lambda,e,D$ satisfying these conditions, there are a
chart family $a$ on the same ambient space and an ambient open set
$N\supset D\cup\operatorname{fr}_X R$ such that $R$ is irreducible in $a$.
The identity on $R$ is chartwise PL from $e$ to $a$ and from $a$ to $e$
on $j^{-1}N$.
\end{theorem}

This requires both PL identity maps near the protected sets, not a global
PL equivalence of structures. The finite-cut argument and its three
specializations are developed next.

\section{A finite bound before choosing a maximal cut}
\label{sec:prime-replacement}

A marked sphere cut consists of finitely many disjoint parametrized PL
spheres in $\operatorname{int}_X R$, disjoint closed product collars with
the spheres at height $1/2$, and the collars' two complete boundary
spheres, called ports. Removing the open collars gives a compact PL
carrier $Q$ whose frontier is the old frontier together with all ports.

Fix a faithful finite PL realization $f:R\to|\mathcal K|$ and its inverse.
A component of $Q$ is a punctured sphere component if, through this
realization, it is finitely PL homeomorphic to a cubical three-sphere with
finitely many disjoint open ball interiors removed, with its entire
frontier corresponding to their boundary spheres. A cut is admissible
when none of its components has this model. The realization and complete
boundary correspondence are part of the condition.

Choose the model from the compact domain before choosing a sphere system.
Let $\mathcal J$ be its old-boundary subcomplex. With triangular faces
$\mathcal K_2$ and boundary vertices $\mathcal J_0$, put
\[
 B_0=4\,\#\mathcal K_2+\#\mathcal J_0,\qquad
 b_1(R)=\dim_{\mathbb F_2}H_1(R;\mathbb F_2).
\]
The implemented uniform bound for an admissible system of $s$ spheres is
\begin{equation}
 s+1\leq b_1(R)+B_0.\label{eq:prime-sphere-bound}
\end{equation}

First normalize the admissible system while retaining its sphere index
set and the fixed original complex. Original chart-star motions avoid the
old boundary and vertices and give finite edge position. Relative motions
remove returning face arcs, and circle surgery removes closed face curves
while retaining the absence of punctured components. Tetrahedral surgery
then strictly decreases the finite total boundary excess. At zero excess,
the centered collars exclude closed interior pieces and supply proper disk
sections and actual ball partitions in each original tetrahedron. All
these replacements retain the number $s$ of spheres. In what follows,
$v$ and $z$ count the components and zero-first-homology components of
this normalized cut.

The incidence graph of the normalized cut has its actual components as
vertices and collars as edges. It is connected; with $v$ vertices its cycle rank is
$s-v+1$. Collapsing components to vertices kills their homology, while
paths inside components supply a graph section up to homotopy. Component
and graph-cycle classes therefore inject together into $H_1(R)$. To check
independence, apply graph collapse to a relation: it detects the cycle
part and annihilates the component part. Once the cycle part is zero,
injectivity from the cut annihilates the remaining relation. If $z$
components have zero first homology, the other $v-z$ components each
contribute at least one dimension. Consequently
\[
 b_1(R)\geq(v-z)+(s-v+1)=s+1-z.
\]

For the geometric bound, consider an original triangular face with $m$
disjoint normal arcs. Their endpoints lie on different edges, so the arcs
fall into $c\leq3$ nonempty corner families. Consecutive arcs in each
ordered family bound a rectangle, giving $m-c$ distinct rectangular
complementary regions. The $m$ proper arcs cut the triangle into $m+1$
regions. Thus at most $c+1\leq4$ regions are exceptional. Each connected
exceptional region lies in one component of the complement of the center
spheres in the original realization. The images of all these face labels
therefore mark at most $4\,\#\mathcal K_2$ global components. Mark also
the components meeting $\mathcal J$: a point in a boundary simplex can
be joined inside that simplex to one of its vertices without meeting a
sphere. These components are the image of the vertex labels
$\mathcal J_0$, so the union of the two sets of marked components has
size at most $B_0$.

Every unmarked component is covered by the regular tetrahedral pieces.
Their rectangular face coordinates glue the pieces as an interval bundle,
with a free involution exchanging the two ends of each fiber. The original
sphere collars identify each connected endpoint surface with a whole
original sphere. Choose a positive-width trim large enough to contain the
corresponding compact cut component; avoidance of the boundary labels puts
this trim in the interior of $R$. If the involution exchanged two distinct
endpoint components, the trim would be a spherical product, and its
contained cut component, with its complete spherical boundary, would have
the forbidden punctured-sphere model. Hence the endpoint component is
invariant. Its connected double cover of the fiber quotient has nontrivial
mod-two monodromy, giving a nonzero first-homology class. Fiber contraction,
the trim coordinates, and the original collars transport that class to the
actual cut component.

Finally, the component map from the cut with open collars removed to the
complement of the center spheres is injective: retract each half-collar to
its endpoint to recover a path in the cut from a path in that complement.
Every zero-first-homology cut component must therefore map into the marked
union above, and distinct cut components have distinct labels there. This
proves $z\leq B_0$ and hence \eqref{eq:prime-sphere-bound}. The face and
vertex labels belong to the original finite realization; subdivisions
used to describe an individual cut do not enlarge this bound.

The whole lattice handle is not a punctured sphere. Since $k>0$, its
lattice basis supplies a circle section and retraction. Every circle-valued
map on a punctured three-sphere kills every loop: extend the map across
the missing balls and contract loops in the whole three-sphere. A
retraction composed with its section would instead preserve the period
loop in the circle, a contradiction. Thus the empty cut is admissible.
Its possible sphere counts form a nonempty bounded set of natural numbers;
choose an actual cut of largest count.

Inserting any additional interior sphere into this cut must produce a
punctured sphere component, or else increase the admissible count. That
component contains a new port; otherwise it was already a component of
the old cut. For lattice handles, the component is disjoint from the old
boundary and its whole frontier consists of complete ports. These
incidence assertions are needed before the capping argument.

\section{Capping and the irreducibility test}
\label{sec:prime-capping}

Refine collars inside a prescribed open neighborhood of the cut spheres.
In a finite realization add one extra coordinate direction for each port
and cone that port in its own direction. The resulting PL ball caps are
pairwise disjoint, and each meets the old carrier in exactly its port.
Enlarge the original domain slightly before realizing it so that the
charts also cover a neighborhood of the old boundary. The capped domain
has exactly that original boundary left.

Suppose an interior PL sphere in this capped domain has no PL ball
filling. Move it off the finitely many cap balls by a relative PL motion
and pull it back to an interior sphere of the original cut carrier.
Absence of a ball is preserved by this motion. Insert the sphere's
collar. Maximality gives a punctured sphere component incident to a new
port. The bounded-region separation argument below rules out its containing
both new ports. Its remaining boundary consists of old ports and just one new
port. Capping the old ports turns this punctured sphere into a ball with
the new port as boundary. The collar identifies this with a filling of
the moved sphere, and the inverse motion fills the original sphere,
a contradiction.

Both the two-new-port exclusion and the full boundary identification are
necessary. Maximal sphere count alone would leave the nonseparating case
unresolved. The separate capped-sphere test performs these steps before
asserting irreducibility.

The cap construction is relative to the original carrier. Each port is given
its own extra coordinate direction, so the cap is a finite PL ball meeting
the carrier exactly in that port; the caps are pairwise disjoint and are
chosen inside the prescribed collar support. If a sphere avoids all cap
balls, its collar can be inserted without changing the old frontier. The new
punctured component has a new port. To exclude both new ports, lift the
inserted sphere through the lattice covering to a locally flat sphere in
$\mathbb R^3$ and take its bounded region. Distinct deck translates have
disjoint frontiers. Their closed bounded regions are either disjoint or
nested; nesting by a nonzero translation is impossible, since a linear
functional nonzero on that translation attains a maximum on the compact
closure. Thus the entire closed region descends injectively, with its
frontier equal to the original sphere. The product-coordinate bounds on
the lifted frontier also bound its enclosed region, so the descended
region lies in the handle. It is regular closed, and the actual collar has
endpoint points in both its interior and the interior of its closed
exterior. A connected cut component avoiding the middle sphere cannot
contain both such points, and therefore cannot contain both complete new
ports. With exactly one new port, capping every old port fills the component
by a ball, and the collar
transports that filling back across the relative motion. The retained
component and its residual cap regions then define the atlas on the same
ambient space; the old charts are kept on their open complement, while the
cap charts are finite PL on the protected side. Hence both identity maps are
PL on one open neighborhood of the protected ball and the entire old
boundary.

The retained-component construction then returns an atlas to the same
original handle. Its residual closed ball regions lie in the interior;
on the open complement of their finite union the old charts are retained.
Transport checks both directions of the PL identity there. When the cut
spheres avoid a positive-index protected ball $D$, choose the collars in
its complement. The residual-region argument puts $D$ and the entire old
boundary on the retained side, supplying the required $N$.

\section{The three lower-index constructions}
\label{sec:prime-index-cases}

For $i=0$ the domain is the whole three-torus and has no old boundary.
Construct a boundary-relative irreducible atlas with a nonempty open
agreement region $N_0$. A PL ambient motion $F$ in the original structure
moves all of $D$ into $N_0$. Pull the new atlas back by $F$. Irreducibility
is preserved: push a sphere forward, fill it, and pull the ball back.
The agreement neighborhood is $F^{-1}N_0$; composition with $F$ and
$F^{-1}$ gives both PL identity directions. This is why an arbitrary
interior marked ball is allowed in index zero.

For $i=1$ first use the actual closed exterior $E=\overline{R\setminus D}$.
Its maximal cut, capped factors, and retained component produce an atlas
in which $E$ is irreducible, the same $D$ is still a marked PL ball, and
$R$ remains a PL domain. Both identities are PL near
$D\cup\operatorname{fr}_X R$. Reattachment is a separate minimization
argument. Suppose a sphere in $\operatorname{int}_X R$ has no PL ball
filling, and choose such a sphere $S$ minimizing its number of contact
circles with $\operatorname{fr}_X E$, among transverse finite polygonal
positions. The family being minimized includes the physical crossing
charts and the finite degree-two contact graph. Merely reducing a
picture's circle count would not establish admissibility of a successor.

An essential proper disk in $E$ is next chosen with minimal intersection,
and circle elimination gives a circle-free representative. Select an
outermost bigon: its boundary is an arc on $S$ and an arc on the lateral
attaching annulus, sharing precisely their endpoints. If these endpoints
belonged to the same contact circle, the terminal disk union and an
inward relative compression would contradict the disk minimum. Hence
the bigon spans two distinct contact circles.

The marked product neighborhood of this same bigon constructs an inside
ball and a protected cap. Their union gives a replacement sphere $S'$
that still does not bound a PL ball. Let $r:[0,1]\times[-1,1]\to
\operatorname{fr}_X E$ be the embedded PL rectangle supplied by this
construction, and set
\[
 E_r=r(\{0,1\}\times[-1,1]),\quad
 H_r=r([0,1]\times\{-1,1\}),\quad
 C_r'=r(\{0,1\}\times\{-1,1\}).
\]
With $\Gamma=S\cap\operatorname{fr}_X E$, the exact contact formula is
\[
 S'\cap\operatorname{fr}_X E
   =\bigl(\Gamma\setminus(E_r\setminus C_r')\bigr)\cup H_r.
\]
The rectangle interior has no other contacts. Since its end arcs belong
to different circles, deleting those arcs and inserting the two side
arcs joins those circles into one, leaving all other components unchanged.
The contact count strictly decreases. Crossing charts persist off the
joined ball; inside it, the product's half-height side boxes and the
matched corner sectors give the new transverse charts. Thus $S'$ belongs
to the original admissible family, contradicting the sphere minimum.
This establishes irreducibility of the whole handle in the retained atlas.

The equal-label bigon case is excluded by the terminal disk minimum, not by
an appeal to a picture: the bigon and the terminal disk form a new proper
disk with strictly fewer intersection arcs, contradicting minimality. In the
spanning case the rectangle has two distinct contact-circle labels, so the
formula above joins exactly those two circles. The product charts on the two
half-height sides and the corner sectors give the crossing charts for the
replacement sphere, which is why the contact minimum is taken over physical
chart data as well as over the number of circles.

For $i=2$ begin with a maximal cut of $R$. The regular cocore construction
and finite contact reduction produce a cut disjoint from a selected
physical cocore, preserving the same sphere index set and the absence of
punctured components. A PL motion through the protected product moves the
whole family off $D$. Transport the collars and ports under that same
motion. The new cut is still maximal because its cardinality is unchanged
and the bound applies to every admissible cut. Protected capping and atlas
transport now finish the construction. This proves the three cases of
Theorem~\ref{thm:protected-prime-replacement}.

In particular, the cocore reduction preserves the no-punctured-component
predicate and the maximal cardinality; it is not merely an isotopy of an
abstract sphere system. The same physical motion carries every collar and
port, so the protected ball is avoided before the retained atlas is assembled.

\section{Use in the lower-handle argument}
\label{sec:compact-core-contracts}

The lower-handle consumer uses three fixed period-lattice models. In index
zero its Wall input ranges over chart families on the punctured zero-handle
ambient space. In indices one and two it ranges over every open subspace
of the corresponding ambient space and every chart set there. A Wall
theorem only for the compact handle itself would therefore be insufficient.
Theorem~\ref{thm:protected-wall-core} supplies these instances by its
ambient generality.

The prime input ranges over all chart sets on each fixed ambient space
and supplies the protected replacement above. The generalized Dehn input
contains the protected annulus or disk-pair construction and the separate
proper-disk theorem with prescribed boundary values. In indices one and
two, its enclosing-region conclusion concerns the same supplied surfaces:
write $T$ for the annulus or the union of the two disks. It returns a
compact region $P$ and a PL parametrization of its entire frontier, with
\[
 \operatorname{fr}_X P=T\cup A_{3/2},\qquad
 C_1\subset P\subset C_2,\qquad
 j^{-1}C_1\subset\operatorname{int}_R j^{-1}P,
\]
\[
 P\subset\pi(B^I\times\{y:\|y\|_\infty<2\}),\qquad
 P\cap\operatorname{fr}_X R=A_{3/2}.
\]
This is the marked side furnished by
Theorem~\ref{thm:dehn-enclosing-region}. The retained original block chart
contains all of $C_2$, hence all of $P$. The marked zero belongs to both
$\operatorname{int}_X R$ and the relative interior of $P$, so $P$ has an
ambient interior point. In this one chart, the PL Alexander theorem
identifies the compact region with its specified sphere frontier as the
bounded PL ball. Transporting that parametrization back gives a chartwise
PL ball on this same $P$, with boundary exactly $\operatorname{fr}_X P$.
Prime replacement is therefore applied to $D=P$, retaining all the
displayed position conditions and both relative PL identity maps.

Relative torus rigidity is supplied by the relative-rigidity chapter.
Together these give straightening for the empty coordinate index set and for every coordinate
subset of size one or two. The smoothing chapter handles the subsequent
index-three and global triangulation steps.

The construction is guided by Hamilton's compact-core and lower-handle
argument~\cite[Lemma 2, pp.~64--65; pp.~66--68]{Hamilton1976}. The
selected-end filling, original-carrier compression, marked cuts, and
protected atlas transport follow the implemented project route. Internal
derivation numbers are traceability records, not sections of Hamilton's
paper. The proper-disk input and the same-surface enclosing regions are
proved in Chapter~\ref{ch:dehn-surfaces}; relative rigidity is proved in
Chapter~\ref{ch:relative-rigidity}.
